\section{Causally Grounded Generative Student Modeling}
\label{sec:method}

\subsection{Problem formulation}

Let $x_t=(q_t,r_t,y_t,k_t,m_t)$ denote an observed interaction at time $t$:
problem content $q_t$, student response $r_t$, correctness $y_t$, available
knowledge-component metadata $k_t$, and optional interaction metadata $m_t$.
Fields unavailable in a source are omitted rather than imputed.  The history
before a target is $H_t=(x_1,\ldots,x_{t-1})$.  Given a task instruction
$\tau$ and a target query $\tilde q_t$, we predict a structured textual output
$o_t$.

Standard generative prompting estimates
\begin{equation}
    p_\theta(o_t\mid \tau,H_t,\tilde q_t),
    \label{eq:full-context}
\end{equation}
which permits the decoder to retrieve arbitrary history tokens and does not
identify a student state.  We instead factor prediction through a fixed-size
continuous state:
\begin{align}
    z_t &= U_\theta(z_{t-1},x_{t-1}),
    \label{eq:update}\\
    p_\theta(o_t\mid \tau,z_t,\tilde q_t)
        &= G_\theta(\tau,z_t,\tilde q_t).
    \label{eq:decode}
\end{align}
The updater $U_\theta$ and decoder $G_\theta$ share the same language-model
backbone and parameter-efficient adapters.  The state
$z_t\in\mathbb{R}^{L\times S\times d}$ retains a layerwise trajectory with
$S=8$ slots, all addressed by one repeated \studentstate{} token type.  It is a
per-student, time-indexed activation, not a learned student-ID embedding.

The tasks differ only in when the state is read and which output fields are
requested.  T1 and T2 read the state before the target interaction; T3 first
updates the state with the observed wrong response and then requests a
misconception; T4 reads the current state but queries a problem $\Delta$
interactions in the future; and strict T5 updates with the current wrong
response before predicting its recurrence on a later exposure.  T6 predicts a
future-KC-performance band and is reported only as an appendix proxy, never as
a direct latent-mastery label.

\subsection{A recurrent state carried by eight slots}

We implement the updater with two banks of eight occurrences of the same
special token type.  For
an update, the model receives
\begin{equation}
 [\studentstate{}_{\mathrm{in}}^{1:S};\;\operatorname{serialize}(x_t);\;
  \studentstate{}_{\mathrm{out}}^{1:S}].
\end{equation}
All positions use the same tokenizer ID.  At the input positions, we replace
the ordinary token embeddings by projected and normalized copies of the
previous state,
\begin{equation}
    e_{\mathrm{in}}^{1:S}=P_{\mathrm{in}}\operatorname{LN}(z_t).
\end{equation}
The causal hidden activations at the output positions can attend to the
previous state and the new interaction; their layerwise trajectory becomes
\begin{equation}
    z_{t+1}=P_{\mathrm{out}}h(\studentstate{}_{\mathrm{out}}^{1:S}).
\end{equation}
For the first observed interaction, $z_0$ is initialized by a learned base
state associated with \studentstate{}.  The small projections $P_{\mathrm{in}}$ and
$P_{\mathrm{out}}$, layer-normalization parameters, and language-model adapters
are trained jointly.

At prediction time, the decoder receives
\begin{equation}
 [\operatorname{instruction}(\tau);\;\studentstate{}^{1:S};\;
  \operatorname{serialize}(\tilde q_t);\;o_t].
\end{equation}
The embeddings at \studentstate{} are replaced by
$P_{\mathrm{in}}\operatorname{LN}(z_t)$.  The raw history is not included, so
every history-dependent prediction must pass through the state.
This removes the history bypass that weakens one-shot bottleneck masks.  It also
makes inference incremental: a serving system stores one vector per student,
updates it after each interaction, and reuses it for any task query.

We unroll updates over chronological blocks and truncate gradients every $B$
interactions.  The state value is preserved across block boundaries while its
gradient is detached.  This bounds memory independently of the full student
history and permits training the 32B backbone with low-rank adapters.

\subsection{Multi-task structured generation}

All targets use a small structured output language.  Binary outcomes use
fields such as \texttt{<correctness>}; response tasks use
\texttt{<predicted\_response>}; misconception diagnosis uses
\texttt{<misconception>}; and prospective tasks combine multiple fields.  The
generative objective is the target-token negative log likelihood
\begin{equation}
    \mathcal L_{\mathrm{gen}}
    =-\sum_{j=1}^{|o_t|}\log
      p_\theta(o_{t,j}\mid o_{t,<j},\tau,z_t,\tilde q_t).
    \label{eq:gen-loss}
\end{equation}
Instruction, state, and query tokens are masked from the loss.  Task and source
samplers are separated so that a large correctness dataset does not dominate
open-response or misconception supervision.

This schema unifies the interface, not the construct or metric.  Closed-choice
responses remain distinct from open mathematical responses; Eedi diagnosis is
evaluated against its full taxonomy or by open generation; and strict
wrong-response recurrence is not merged with repeat-incorrectness proxies.

\subsection{Predictive state distillation}

A strict bottleneck can lose information before the updater learns useful
compression.  We therefore initialize a full-context teacher $T$ from the
strong unrestricted multi-task model.  For the same target, the teacher sees
$H_t$ as in Equation~\ref{eq:full-context}, while the student decoder sees only
$z_t$.  We minimize a temperature-scaled token-level divergence
\begin{equation}
    \mathcal L_{\mathrm{KD}}
    =\mathcal T^2\sum_j
      \operatorname{KL}\!\left(
      p_T^{(\mathcal T)}(\cdot\mid H_t,\tau,\tilde q_t,o_{t,<j})
      \;\|\;
      p_\theta^{(\mathcal T)}(\cdot\mid z_t,\tau,\tilde q_t,o_{t,<j})
      \right).
    \label{eq:kd-loss}
\end{equation}
The teacher is frozen.  Ground-truth generation remains the primary objective,
so the student is not forced to reproduce teacher errors.  The gap between the
teacher and state-only student is also reported as an empirical measure of
state sufficiency.

\subsection{Recurrent--direct state consistency}

Errors can accumulate when the state is updated over a long trajectory.  We
construct a second state from the complete visible history using a direct
history encoder $E_\theta$ with the existing late-fusion mask.  Its state token
can attend to all of $H_t$, but the answer cannot.  For selected checkpoints in
an unrolled trajectory, we align the recurrent state $z_t^{\mathrm{rec}}$ with
the stop-gradient direct state $z_t^{\mathrm{dir}}=E_\theta(H_t)$:
\begin{equation}
    \mathcal L_{\mathrm{dyn}}
      =1-\cos\!\left(
        \operatorname{LN}(z_t^{\mathrm{rec}}),
        \operatorname{sg}[\operatorname{LN}(z_t^{\mathrm{dir}})]
      \right).
    \label{eq:dyn-loss}
\end{equation}
We reverse the stop-gradient direction on half of the minibatches.  This
symmetric schedule prevents one branch from remaining an untrained target
while avoiding immediate representational collapse.  Both states must also
predict the same observed target through Equation~\ref{eq:gen-loss}; cosine
agreement alone is never optimized without behavioral supervision.

\subsection{Counterfactual grounding with executable misconceptions}

Decodability does not ensure that changing a state changes behavior.  We use
MalruleLib to construct supervised interventions.  Let histories $H^A$ and
$H^B$ instantiate different misconception rules $a$ and $b$, and let $q$ be a
held-out problem for which both rules are executable.  The library provides
counterfactual responses $r_a(q)$ and $r_b(q)$.  We decode the same query after
swapping the two states and optimize
\begin{align}
    \mathcal L_{\mathrm{cf}} ={}&
    -\log p_\theta(r_a(q)\mid z^A,\tau_{\mathrm{response}},q) \\
    &-\log p_\theta(r_b(q)\mid z^B,\tau_{\mathrm{response}},q).
    \label{eq:cf-loss}
\end{align}
The pairing procedure deliberately changes the problem template, so the model
cannot satisfy the loss by copying a previous answer.  We do not attach a
malrule classifier to $z$: the supervision specifies what the state should
\emph{do}, not which label a probe should recover.  Evaluation uses unseen
history--query pairs, held-out templates, and, where coverage permits, held-out
malrules.

\subsection{Future-profile ownership grounding}

To make student identity testable without exposing a student ID, we add an
auxiliary readout over a fixed projection $\phi(z_t)$ of the layerwise State8
trajectory.  For student $i$, let $y_i\in[0,1]^K$ be smoothed correctness rates
computed only from post-cutoff interactions and let $m_i$ mark KCs with
sufficient support.  A linear head predicts
\begin{equation}
    \hat y_i=\sigma(W\phi(z_i)+b).
\end{equation}
For a hard mismatched student $j$, selected inside the training split, we
optimize supported future-rate error together with an ownership margin:
\begin{align}
    \mathcal L_{\mathrm{prof}} &=
      \operatorname{MSE}_{m_i}(\hat y_i,y_i),\\
    \mathcal L_{\mathrm{own}} &=
      \max\{0,\gamma+
      \operatorname{MSE}_{m_i\cap m_j}(\hat y_i,y_i)-
      \operatorname{MSE}_{m_i\cap m_j}(\hat y_i,y_j)\}.
\end{align}
Post-cutoff targets supervise training but are never visible to the updater or
decoder at inference.  KCs below a threshold fixed from training support are
excluded from this auxiliary loss.

The complete objective is
\begin{equation}
    \mathcal L
      =\mathcal L_{\mathrm{gen}}
      +\lambda_{\mathrm{KD}}\mathcal L_{\mathrm{KD}}
      +\lambda_{\mathrm{dyn}}\mathcal L_{\mathrm{dyn}}
      +\lambda_{\mathrm{cf}}\mathcal L_{\mathrm{cf}}
      +\lambda_{\mathrm{prof}}\mathcal L_{\mathrm{prof}}
      +\lambda_{\mathrm{own}}\mathcal L_{\mathrm{own}}.
    \label{eq:total-loss}
\end{equation}
We first warm-start on unrestricted multi-task generation, then train recurrent
and direct state paths with distillation and consistency, and finally mix
counterfactual and profile batches with real-data replay.  This curriculum separates
language/output-format learning from state compression and reduces catastrophic
forgetting during causal grounding.
